\subsection{Origin of the same-event background in $m^2$}
\label{sec:bkg}

Sec.~\ref{sec:m2} found that the overlay mixed event does not describe the same-event (SE) $m^2$ away from the peaks.
In particular, the SE yield between the proton and deuteron peaks ($1.6<m^2<2.8$\,GeV$^2$/$c^4$, called HI below)
is several times larger than the mixed-event (ME) estimate. This section reports the study of that mismatch.
It uses four extra passes of \fpath{StRoot/StTofMEMaker} over the same 15 runs and events as Sec.~\ref{sec:me1m}:
\begin{itemize}
\item v4: SE and ME split by 12 track, event and hit variables;
\item v5: the same quantities per projected BTOF cell;
\item v6: details of neighbour-cell matches, on the first 25\,000 events of each run;
\item v7, v8: two emulations of the production hit choice.
\end{itemize}
Scripts: \fpath{scripts/make_bkg.C}, \fpath{make_cell.C}, \fpath{make_nb.C}, \fpath{make_m2.C}, \fpath{make_late.C};
logs \fpath{logs/bkg_v4.txt}, \fpath{cell_v5.txt}, \fpath{nb_v6.txt}, \fpath{m2_numbers_v7.txt}, \fpath{m2_numbers_v8.txt},
\fpath{late_v8.txt}. All comparisons are for 0--80\% centrality unless stated otherwise.

\subsubsection{Where the HI excess is}
Table~\ref{tab:bkgcat} lists the HI yield per $10^4$ projected tracks for SE and for the overlay ME of
Sec.~\ref{sec:m2}, split by each variable.
\begin{itemize}
\item \textbf{No dependence:} the SE/ME ratio does not depend on the difference between the TPC and VPD vertex
  ($\approx6$ in every class, including events without a VPD vertex), on $|z_{\rm local}|$, on $\eta$ or on the track
  charge. This rules out start-time errors and event pile-up as the main source
  (Fig.~\ref{fig:bkg_v4_k1_p08}, \ref{fig:bkg_v4_k11_p08}).
\item \textbf{Mild dependence:} the ratio rises with the DCA ($5.9\to9.0$) and is larger for short tracks.
  Tracks with $n\sigma_{\rm p}>8$ (light-nucleus candidates) show a strong excess, but they are a small share.
\item \textbf{Strong dependence, the chosen cell:} the HI excess comes from tracks matched to a hit in a
  \emph{neighbouring} cell (Fig.~\ref{fig:bkg_v4_k7_p08}). Neighbour-cell matches are 6\% of the SE matches but
  \nbShareHI\% of the SE HI entries (v6). SE/ME in the HI window is 10.5 for them and 2.0 for projected-cell matches.
\end{itemize}

The neighbour-cell excess is not caused by dead or inefficient cells. Cells with an efficiency above 0.8 (79\% of the
projected tracks) carry \clShared\%+\clSharee\% of the HI entries, and cells below 0.5 carry about 1\%
(Table~\ref{tab:celleff}, Fig.~\ref{fig:cell_v5_eff}). It is not a projection effect either.
The rate of HI neighbour matches is about $7\times10^{-4}$ per SE match whether the track crosses the centre or the
edge of the projected cell (Fig.~\ref{fig:nb_v6}). At the centre the neighbour cell is $\ge1.7$\,cm away, against a
projection precision of about 0.05\,cm at these momenta. Finally, in \nbStUntN{} of the \nbStTotN{}
neighbour matches (v6, $0.6<p<1.5$\,\GeVc) the projected cell \emph{did} contain a hit, but that hit has no time in
the picoDst: production did not match it to any track and gave our track the neighbour hit instead.
The HI fraction of these matches is \nbStUntHIPK{} per mille of the peak yield, against \nbStEmptyHIPK{} per mille when
the projected cell was empty.

\begin{table}[H]
\centering
\caption{HI yield ($1.6<m^2<2.8$\,GeV$^2$/$c^4$) per $10^4$ projected tracks for SE and overlay ME (projected cell
first, all hits), by category. ``share'' is the category's share of all SE matches.
$^{*}$Hit variables: absolute counts instead of rates, because projected tracks cannot be assigned to them.
Generated: \fpath{tables/bkg_cat_v4.tex}.}
\label{tab:bkgcat}
\scriptsize
\input{tables/bkg_cat_v4.tex}
\end{table}

\begin{table}[H]
\centering
\caption{HI yield by cell efficiency $\varepsilon_c$ (fraction of tracks projected to the cell for which the cell
has a hit), $0.6<p<3$\,\GeVc. ``nb.\ match'': share of the projected tracks matched to a neighbour cell.
Generated: \fpath{tables/cell_eff_v5.tex}.}
\label{tab:celleff}
\small
\input{tables/cell_eff_v5.tex}
\end{table}

\subsubsection{Production chooses among several hits by ToT}
In \fcmd{StBTofMatchMaker} (SL24y, lines 782--833) a track with hits in more than one candidate cell (projected cell
and neighbours) is matched to the hit with the \emph{largest ToT} (match flag 2). Only for equal ToT is the hit
closest in $y$ used (flag 3). So a random hit in a neighbour cell can take the place of the track's own hit. The overlay
ME of Sec.~\ref{sec:m2} always kept the own hit in the projected cell and therefore misses these cases. The picoDst
has no ToT, so two emulations were run:
\begin{itemize}
\item \textbf{Equal chance} (v7): ToT independent of the hit's origin. Each of the hits in the candidate cells,
  from this event and from the partner, wins with probability $1/(n_A+n_B)$. Untimed partner hits are included
  through the scaling of Sec.~\ref{sec:m2}.
\item \textbf{Earliest time wins} (v8; ToT replaced by the time, smaller time wins): the hit chosen for the track in
  this event competes with the timed partner hits in the candidate cells. Untimed partner hits cannot take part.
\end{itemize}

\begin{center}
\small
\begin{tabular}{lccc}
\toprule
 & proj.\ cell first & equal chance & earliest time \\
\midrule
random matches / SE matches, 0--80\% (\%) & \mtMENormFracVthree & \mtMENormFracVseven & \mtMENormFracVeight \\
\quad 0--5\% (\%) & \mtMENormFracCiVthree & \mtMENormFracCiVseven & \mtMENormFracCiVeight \\
\quad 70--80\% (\%) & \mtMENormFracCaVthree & \mtMENormFracCaVseven & \mtMENormFracCaVeight \\
max ME/SE per $m^2$ bin, 0--5\%, $0.4<p<1.5$ & \mtValleyAllVthree & \mtValleyAllVseven & \mtValleyAllVeight \\
SE/ME, HI window, 0--80\%, $0.2<p<3$ & \mtSbRatioAllVthree & \mtSbRatioAllVseven & \mtSbRatioAllVeight \\
\bottomrule
\end{tabular}
\end{center}
Under the peaks in 0--5\% at $0.8<p<1$\,\GeVc{} the random fraction is \mtFrPiMidCiVseven\% ($\pi$), \mtFrKMidCiVseven\% (K) and
\mtFrPMidCiVseven\% (p) with equal chance, and \mtFrPiMidCiVeight\%, \mtFrKMidCiVeight\% and \mtFrPMidCiVeight\% with the earliest-time
rule (Figs.~\ref{fig:m2_v7_cent8}, \ref{fig:m2_v8_cent8}).
\begin{itemize}
\item Equal chance roughly doubles the random matches. It describes the HI region above $p\approx0.8$\,\GeVc, but it
  exceeds the SE yield in the valleys next to the peaks by up to a factor \mtValleyAllVseven, which is impossible.
  The track's own hit must therefore win more often than an equal-chance draw gives.
\item Earliest time wins never exceeds the SE yield (largest ratio \mtValleyAllVeight). By construction, though, a random
  hit wins only when it is earlier than the own hit, so it fills the low-$m^2$ side and leaves the HI excess
  (SE/ME = \mtSbRatioAllVeight).
\item Production selects on ToT, not on time, so the two emulations bracket the in-time random matches that come from
  hits with a time: \mtMENormFracCiVeight--\mtMENormFracCiVseven\% of the TOF matches in 0--5\%.
\end{itemize}

\subsubsection{A late component flat in time}
In time units the HI excess is a nearly flat plateau. Between 16 and 26\,ns after the pion arrival time the SE rate
falls only slowly, while the ME falls by about a factor 4. Only 0.45\% of the hits that have a time in the picoDst arrive
later than 25\,ns. The plateau hits must therefore be among the hits without a time, i.e.\ hits that production did
not match to any track in their own event. Such a hit enters the SE when production matches it to the track under
study, because the track's BTofPidTraits then carries its time. In a partner event it has no time and cannot give an
$m^2$, so no picoDst-based mixed event can contain it.

The plateau rate per projected track (Fig.~\ref{fig:late_v8}) rises only from \lateHiPer{} to \lateHiCen{} hits per ns
per $10^6$ tracks (factor \lateHiRatio) from 70--80\% to 0--5\% for $0.6<p<1.5$\,\GeVc, and from \lateLoPer{} to \lateLoCen{}
(factor \lateLoRatio) for $0.2<p<0.6$\,\GeVc. The BTOF hit multiplicity, which follows refMult (Fig.~\ref{fig:evt_refmult_vs_tofmult}), grows by more
than an order of magnitude over the same range (read from the refMult range of the classes, not computed per class).
Most of the plateau is therefore unrelated to the triggered collision: electronic noise or
hits from other bunch crossings in the TOF readout window.

\caveat{The picoDst cannot provide a complete mixed-event description of random matches: (1) it has no ToT, so the
production choice among several hits can only be bracketed; (2) 65\% of the BTOF hits have no time, among them the flat
late component. Both are available in the MuDst (\fcmd{StMuBTofHit}), where every hit has a time and a ToT; a MuDst
test is being set up (Sec.~\ref{sec:open}).}

\begin{figure}[H]
\centering
\figpanel{0.98}{bkg_v4_k7_p08}{SE (points) and overlay ME (red) $m^2$, $0.8<p<1.5$\,\GeVc, 0--80\%, for hits in the projected
cell (left) and in a neighbour cell (right). Per $10^4$ projected tracks is not defined for hit categories; ME is normalized
per projected track in total.}
\\[1ex]
\figpanel{0.98}{bkg_v4_k1_p08}{Same, split by $|V_z-V_z^{\rm VPD}|$: $<3$, 3--6, $>6$\,cm and no VPD vertex. ME normalized per
projected track of each class.}
\\[1ex]
\figpanel{0.49}{bkg_v4_k11_p08}{Same, split by charge.}
\hfill
\figpanel{0.49}{cell_v5_eff}{Distribution of the cell efficiency $\varepsilon_c$ ($0.6<p<3$\,\GeVc), log scale.}
\caption{Background diagnostics (v4, v5).}
\end{figure}

\begin{figure}[H]
\centering
\figpanel{0.98}{nb_v6}{Left: HI entries per $10^4$ SE matches vs $|y_{\rm local}|$ of the projection for SE (black) and overlay
ME (red), and the share of neighbour-cell matches (open blue, \%). Right: production $y_{\rm local}$ of own neighbour hits
(read by eye; the quantity is not used in the conclusions).}
\\[1ex]
\figpanel{0.6}{late_v8}{SE hits arriving 16--26\,ns after the pion time, per ns and per $10^6$ projected tracks, vs centrality
class, for $0.2<p<0.6$ (blue) and $0.6<p<1.5$\,\GeVc{} (red).}
\caption{Neighbour-cell matches (v6) and the late plateau (v8).}
\end{figure}

\begin{figure}[H]
\centering
\figpanel{0.82}{m2_v7_cent8}{$m^2$ of projected QA tracks, 0--5\%, nine momentum bins (log scale). Black: same event.
Red filled: overlay ME with the equal-chance production choice, all hits, normalized per projected track. Blue dashed: overlay
ME with the projected cell first (Sec.~\ref{sec:m2}).}
\caption{SE vs ME $m^2$, 0--5\%, equal chance.}
\end{figure}

\begin{figure}[H]
\centering
\figpanel{0.82}{m2_v8_cent8}{$m^2$ of projected QA tracks, 0--5\%, nine momentum bins (log scale). Black: same event.
Red filled: overlay ME with the earliest time winning, timed partner hits only, normalized per projected track. Blue dashed:
equal-chance choice.}
\caption{SE vs ME $m^2$, 0--5\%, earliest time wins.}
\end{figure}
